\scanpage{181}
\mspage{37}

3. \uline{Retour sur la cofinalité}.

Il est instructif de relier les questions de cofinalité (pour une flèche $u:I'\mathbin{-}I$ dans $\mathrm{Cat}$) au dérivateur $\mathrm{HOT}_{W^{\deleted{\unsure{*|\infty}}}_0}$, où $W^{\deleted{\unsure{*|\infty}}}_0\subset\mathrm{Fl}\,\mathrm{Cat}$ est le plus grossier des localiseurs dans $\mathrm{Cat}$ satisfaisant l'axiome de connexité $\mathrm{Loc}\,5$,

\unsure{i.e}

\[
W_0=\bigl\{f\in\mathrm{Fl}\,\mathrm{Cat}\ \big|\ \pi_0(f)\ \text{bijectif}\bigr\}\tag*{(3.1.)}
\]

On trouve alors, \unsure{pour} $I$ dans $\mathrm{Cat}$, une équivalence canonique, $2$-fonctorielle en $I$

\[
\mathrm{HOT}_{W_0}(I)\xrightarrow{\approx}I^\wedge\overset{\mathrm{def}}{=}\underline{\mathrm{Hom}}(I,(\mathrm{Ens}))\tag*{(3.2.)}
\]

en d'autres termes on a une équivalence

\[
\mathrm{HOT}_{W_0}\simeq\mathbb D_{(\mathrm{Ens})}\,.\tag*{(3.3.)}
\]

Quand $\mathrm{HOT}_{W_0}(I)$ est défini par localisation à partir de la catégorie $\mathrm{Cat}(I)\overset{\mathrm{def}}{=}\underline{\mathrm{Hom}}(I^o,\mathrm{Cat})$, la flèche 3.2. est celle déduite, par localisation, de la flèche

\[
\begin{aligned}
\pi_0^I:\ \underline{\mathrm{Hom}}(I^o,\mathrm{Cat})&\unsure{\longrightarrow|\mathbin{-}}\underline{\mathrm{Hom}}(I^o,(\mathrm{Ens}))\\
X\deleted{\unsure{_{*}|_{/}}}&\longmapsto\pi_0\circ X
\end{aligned}\tag*{(3.4)}
\]

\deleted{\unsure{$\pi$}}

\[
\pi_0:\mathrm{Cat}\mathbin{-}(\mathrm{Ens})\tag*{\unsure{(3.5)|(2.5)}}
\]

est le foncteur ``composantes connexes''.

Si l'on décrit $\mathrm{HOT}_{W_{\unsure{0|1}}}(I)$ comme localisé de $\mathrm{Cat}_{/I}$, (3.2), est déduite de la flèche

\[
\left\{\begin{aligned}
\mathrm{Cat}_{/I}&\longrightarrow\deleted{\unsure{I^\wedge|I}}\,I^\wedge=\underline{\mathrm{Hom}}(I^o,(\mathrm{Ens}))\\
X&\longmapsto\deleted{\unsure{\xi|X}}\,\pi_0(X/I)=(i\mapsto\pi_0(i\backslash X))\end{aligned}\right.\tag*{(3.6)}
\]

% archive sheet number: 181
